\section{数据集、训练与实验计划}
\label{sec:experiment}
\subsection{数据集}
拟采用 RAF-DB、KMU-FED、PPB-Emo 和 FER2013/FER+ 构建训练与评价数据。选择依据是通用表情覆盖与近红外座舱成像互补：通用数据提供多样的表情监督，座舱数据用于确认目标场景中的动作与纹理。
\begin{table}[htbp]
\centering
\caption{数据集选择与拟承担任务}
\label{tab:datasets}
\begin{tabularx}{\linewidth}{@{}lXX@{}}
\toprule
数据来源 & 数据特点 & 拟承担任务\\
\midrule
RAF-DB & 七类自然场景表情 & 通用特征学习与验证\\
KMU-FED & 真实车辆近红外面部图像 & 座舱适配和匿名化对照\\
PPB-Emo & 驾驶任务中央红外视频 & 补充中性与视频样本\\
FER2013/FER+ & 图像与多标注者投票 & 筛选一致性较高的补充样本\\
\bottomrule
\end{tabularx}
\end{table}
数据来源分别见 RAF-DB、KMU-FED、PPB-Emo 和 FER+ 的公开论文\cite{li2017rafdb,kmu2018,ppbemo,barsoum2016ferplus}。整理时建立统一七类标签映射、来源记录与许可说明，按实际用途取样。

\subsection{预处理}
拟统一灰度化、人脸检测、五点对齐与尺寸调整，并在匿名化后重新提取分类输入。先确定训练、验证和测试用途，再生成匿名缓存，记录源文件与处理配置。RAF-DB 沿用官方划分，KMU-FED 保留序列尾帧作回归，PPB-Emo 按参与者隔离划分，保证评价条件清楚。

\subsection{训练计划}
先训练轻量卷积表情基线，确认标签、归一化和评价程序正确；再加入车载样本，通过加权采样平衡数据来源；最后在匿名缓存上微调，使模型直接学习匿名域的表情结构。拟使用交叉熵监督真实标签，并比较轻微亮度、平移等增强及一致性约束的作用。

匿名化训练拟固定身份与表情评价接口，通过调整局部编辑幅度与损失权重形成候选。对每个候选使用同一验证集合，依次比较身份命中、表情准确率和处理时间。质量配置可进一步比较预训练表征、翻转增强与温度校准，只有验证结果支持时才纳入应用。

\subsection{实验指标}
匿名化评价只针对自建身份模型，固定脸部裁剪范围后比较原身份命中与特征相似程度，并核对掩膜外像素。表情评价采用准确率、分类别召回与混淆矩阵；视频评价记录完成分析的帧率、延迟和类别切换情况。各集合独立报告，不合并为一个缺少来源的总分。

\subsection{项目分工}
\begin{table}[htbp]
\centering
\caption{小组工作安排}
\label{tab:division}
\begin{tabularx}{\linewidth}{@{}lX@{}}
\toprule
成员 & 拟负责内容\\
\midrule
高霁 & 总体方案与核心链路，匿名化及表情模块整合、桌面交互、GPU 推理优化，实验统筹、交付整理和报告统稿\\
同学1 & 数据资料与许可整理、样本清洗、灰度化和人脸对齐\\
同学2 & 表情训练方案、参数记录与分类别结果分析\\
同学3 & 匿名化调研、备选方案比较与面部编辑分析\\
同学4 & 视频与摄像头输入、连续推理联调和性能记录\\
同学5 & 对照实验图表、文档校对与演示材料整理\\
\bottomrule
\end{tabularx}
\end{table}

\subsection{进度安排}
项目按需求调研、数据准备、模型训练、桌面集成、联合评价和材料整理推进。数据阶段完成标签与划分确认后再训练，模型阶段完成固定集合的比较后再接入应用，最后统一核对代码、权重与报告中的实验设置。
\begin{table}[htbp]
\centering
\caption{阶段任务与阶段成果}
\label{tab:schedule}
\begin{tabularx}{\linewidth}{@{}lXX@{}}
\toprule
阶段 & 主要工作 & 阶段成果\\
\midrule
调研与准备 & 需求、方法、数据来源比较 & 技术路线与数据清单\\
训练与选择 & 基线、匿名化及匿名域适配 & 候选与分域验证记录\\
系统集成 & 图片、视频、摄像头与 GPU & 可操作桌面应用\\
评价与交付 & 联合评价、运行说明与汇报 & 报告、演示及源码资料\\
\bottomrule
\end{tabularx}
\end{table}
